\section{Methods}
\label{sec:methods}

\subsection{Predictive Models}

Table~\ref{tab:models} lists the five models. All receive the same 33 inputs, with coded
administrative attributes (for example course, application route and parents' occupations) treated
as categorical. Logistic regression, the decision tree and XGBoost share one preprocessing step
(one-hot encoded categoricals; median-imputed, standardised numericals), whereas EBM and \tabicl
receive the named columns directly. All settings were fixed in advance in a written analysis plan
rather than tuned on the test set.

\textbf{Model complexity.} The models differ in how far a person can follow them. \tabicl is a
27.6-million-parameter transformer with 18 attention blocks (3 column, 3 row and 12 in-context),
which scores each student by attending over all 2,541 training students as context and averages 4
ensemble members. It is therefore neither simulatable nor decomposable, since its reasoning cannot be
recovered from a coefficient list or a small tree, and it requires post-hoc explanation. The main-effects EBM, in contrast, is
interpretable by construction: each prediction is exactly the logistic transform of an intercept
plus one learned shape-function contribution per feature (its validation probabilities are
reconstructed from these terms within $3.3\times10^{-16}$). Logistic regression and the depth-3 tree
are also simulatable.

\textbf{Validation-stage selection.} On the validation split, ROC-AUC was 0.952 for \tabicl, 0.946
for EBM, 0.945 for logistic regression and XGBoost, and 0.923 for the tree. As the margin of \tabicl
over EBM (0.006) was modest, the proposal had already identified EBM as a potentially more
defensible candidate and deferred the decision to the final test.

\begin{table}[!htb]
\centering
\caption{Models compared at the same 20\% outreach capacity.}
\label{tab:models}
\footnotesize
\setlength{\tabcolsep}{3pt}
\begin{tabular}{@{}p{2.1cm}p{5.9cm}@{}}
\toprule
\textbf{Model} & \textbf{Configuration} \\
\midrule
Logistic regression & One-hot categoricals, balanced class weights \\
Decision tree & Depth 3, balanced class weights \\
XGBoost & 300 trees, depth 4, learning rate 0.05, subsampling 0.8~\citep{chen2016xgboost} \\
EBM & Main effects only (no interactions), interpret 0.7.8~\citep{lou2013ebm} \\
\tabicl & v2.2.0, 4 ensemble members, CPU, official checkpoint~\citep{qu2025tabicl} \\
\bottomrule
\end{tabular}
\end{table}

\subsection{Predictive and Capacity-Based Evaluation}

Ranking quality is measured by ROC-AUC, and probability quality by the Brier score and the 10-bin
expected calibration error (ECE). Since the decision is to ``flag the top 20\%'', the principal
operating metrics are dropout recall and precision among flagged students:
\begin{align}
\text{Recall@20\%} &= \frac{\#\{\text{flagged dropouts}\}}{\#\{\text{dropouts}\}}, \label{eq:recall}\\
\text{Precision@20\%} &= \frac{\#\{\text{flagged dropouts}\}}{\#\{\text{flagged}\}}. \label{eq:precision}
\end{align}
With a fixed capacity, recall is capped: with 213 dropouts and 109 slots in the test set, at most
51\% of dropouts can be flagged. Uncertainty in the differences between \tabicl and each baseline is
estimated with 2,000 paired bootstrap resamples of the test set (seed 42).

\subsection{Local and Global Explanations}

Local explanations are computed with permutation SHAP~\citep{lundberg2017shap} on the predicted
dropout probability, using 25 training rows as the background (the ``typical student'' reference), kept small since
every SHAP evaluation scores the student against each background row.
Signed SHAP values of related features are summed into semantic groups, such as
\emph{first-semester progress} (units passed, grade and pass rate) and \emph{first-semester
participation} (units enrolled and evaluated), and the five groups with the largest absolute values
are displayed. The SHAP budget is \texttt{max\_evals}$\,=67=2\times 33+1$ model evaluations per
student, which is exactly one forward-and-backward pass over one random feature ordering, the
smallest budget shap's permutation explainer accepts. Exact TreeSHAP is
unavailable, as \tabicl is not a tree ensemble. Before the final test, the pipeline was verified:
SHAP values reconstructed \tabicl's probabilities within $4.1\times10^{-8}$, and on a synthetic
white-box model the explainer recovered the known feature ranking exactly (Spearman 1.0). Global
behaviour is summarised by grouped permutation importance and by ALE~\citep{apley2020ale}, which is
less affected than partial dependence by correlated features.

\subsection{Explanation Audit}

The analysis plan originally scheduled the full audit on 20 validation students. Owing to the project
timeline, it was instead conducted after the final test on all 544 validation students, using a
\tabicl refit with the frozen settings. This refit reproduces the frozen model on validation
(largest probability difference $1.0\times10^{-5}$; all 108 flags identical); it therefore describes
the tested model, although it constitutes post-hoc evidence rather than a gate passed before the
final test.
\begin{itemize}[nosep,leftmargin=*]
  \item \textbf{Faithfulness.} Groups are replaced by training medians or modes in SHAP order and in
  50 random orders; a positive area between the two cumulative probability-change curves indicates
  that the SHAP ranking identifies the groups on which the model relies.
  \item \textbf{Stability.} Each explanation is recomputed 10 times under three variations: a
  different random seed, a different 25-row background sample and minimal grade perturbations
  (Gaussian noise of 1\% of the training standard deviation, clipped to the training range and
  accepted only if the prediction moves by at most 0.01). Stability is the mean pairwise overlap of
  the top-five groups, since five bars are what an adviser sees; values below 0.80, meaning that more
  than one of the five reasons typically changes, are flagged as unstable. All 544 students are used
  so that instability can be counted rather than illustrated.
  \item \textbf{Constrained recourse.} Features were classified in advance by availability and
  feasibility: age, nationality, prior qualifications and completed first-semester results remain
  fixed. For flagged students, only the administrative fields \emph{debtor} and \emph{tuition fees
  up to date} may change, as these could realistically be resolved, often with institutional help,
  before the second semester. The smallest such change is deemed to work if the student leaves the
  top 20\% when re-ranked with all other scores fixed; students without a working change are
  reported explicitly.
  \item \textbf{Surrogates and pairs.} Six local linear surrogates (OLS, ridge and lasso, each with
  and without pair terms) and pairwise group masking assess how additive the model's local behaviour
  is.
\end{itemize}

\subsection{Human-Study Analysis}

The pre-specified primary outcome is each participant's accuracy over the ten cases, defined as
agreement with the model's flag. Group means are reported with 95\% intervals from 10,000 stratified
bootstrap resamples, and the difference between Group~B and Group~A with its interval, a two-sided
permutation $p$-value (10,000 label permutations) and Hedges' $g$. With 25 participants per group,
these comparisons are exploratory. Confidence and self-reported time are analysed in the same way,
while case-level accuracy, a leave-one-case-out sensitivity analysis and the chart-reading checks
are reported descriptively.

\subsection{Qualitative Coding}

The 1,500 written answers (50 participants $\times$ 10 cases $\times$ 3 questions) were coded with a
pre-written codebook, assigning one primary theme per answer before group membership was attached.
Each support suggestion was additionally rated on four items from the proposal: whether it names a
concrete next step, whether it could begin before the second semester, who is responsible (the
student, the institution or both) and whether it asserts an unsupported cause. The planned
independent double-coding of 20\% of answers was not carried out; inter-coder agreement was
therefore not measured, and the themes represent a single descriptive coding pass.
